\documentclass[11pt]{article}
\usepackage[margin=2.2cm]{geometry}
\usepackage{booktabs}
\usepackage{amsmath}
\usepackage{graphicx}
\usepackage{hyperref}

\title{Agent360: A Concealment-First Bilateral Negotiator\\for ANL 2026}
\author{%
  Noam Kazum \\[0.4em]
  Omer Shani Steinmetz \\[0.4em]
  Dr.~Galit Haim \\[0.4em]
  Dr.~Raz Lin \\[1em]
  \textit{The College of Management Academic Studies}%
}
\date{June 15, 2026}

\begin{document}
\maketitle

\begin{abstract}
We present \textbf{Agent360}, a bilateral negotiator submitted to the Automated Negotiation
League (ANL) 2026.
The competition scores each agent on \textbf{Advantage} (deal quality relative to a reserved
value) and \textbf{Concealing} (how poorly the opponent learns our true preferences from our
bid stream).
Our agent design follows a \emph{concealment-first} principle: early offers
deliberately misrepresent issue priorities so that frequency-based opponent learners infer an
incorrect utility function, while later phases and acceptance rules extract profitable
agreements.
The agent uses no oracle access to the opponent's class or true utility--only observed offers,
negotiation seat, and the relative deadline time.
We relate the bidding persona to theories of signaling, strategic misrepresentation, and the
limits of honesty in non-cooperative negotiation.
\end{abstract}

\section{Introduction}

In bilateral alternating-offer negotiation, each proposal is both an economic act and a
\emph{signal} about private preferences.
The ANL 2026 league makes this dual role explicit: an agent is rewarded not only for closing
good deals, but also for preventing the opponent from building an accurate model of its true
utility function.
Formally, the league score is
\begin{equation}
  \text{Score} = \text{Advantage} + \text{Concealing},
\end{equation}
where concealment is derived from the rank correlation between the agent's true utility and the
utility function the opponent inferred from the negotiation transcript.

Most student agents in the league learn from offers using Smith-style frequency models or
similar adaptive estimators.
Against such opponents, bidding one's sincere top outcomes early tends to maximize short-term
deal quality but minimizes concealing, because the opponent's learner quickly converges on the
correct issue weights.
Agent360, therefore, treats preference \emph{concealment} as a first-class design goal and tunes deal extraction only in ways that do not collapse the decoy persona.

This report describes the agent, its psychological motivation, its
three-layer architecture, and the evaluation process that produced the final parameterization.

\section{Background: Deception, Signaling, and Learning}

\subsection{Offers as preference signals}

In multi-issue negotiation, an offer is a vector of issue values.
Rational opponents treat repeated offers as evidence about the values the sender prefers.
Smith-style frequency models count how often each value appears in the opponent's
bids, implementing a simple form of \emph{preference elicitation from behavior}.
If our early bids consistently favor values that are actually low-priority for us, the
opponent's estimated utility function shifts away from the truth, which directly increases concealment under the league metric.

This is distinct from verbal deception: ANL agents cannot negotiate in natural language.
All misdirection must be encoded in the \emph{bid stream}, analogous to non-verbal signaling
in human bargaining (consistent choices that suggest false priorities).

\subsection{Truth mechanisms versus competitive misrepresentation}

Mechanism-design results (e.g.\ strategy-proof auction formats) show that honesty can be a
dominant strategy when the \emph{rules} reward truthful reporting.
The ANL scoring rule has the opposite incentive on the concealment axis: if the opponent models the agent correctly, our concealment score falls.
We therefore operate in a non-cooperative setting closer to competitive signaling than to
 incentive-compatible mechanism design.

We initially considered a \emph{reverse-psychology} persona: bid sincere high-utility outcomes early (as if ``playing truth''), then shift toward misdirection, exploiting opponents that systematically discount our offers.
Local experiments showed that this variant improved advantage against some time-based
negotiators but reduced concealment against BOA, MAP, and MiCRO-style learners---the dominant
opponent family in both our benchmarks and the live tournament.
We rejected reverse-truth as the submission strategy and retained a gradual decoy-to-truth
persona instead.

\subsection{Opponent modeling without labels}

Human negotiators form \emph{theories of mind} about whether the partner is conceding,
posturing, or mirroring prior bids.
Agent360 approximates this with a lightweight behavioral classifier (\texttt{mirror},
\texttt{learner}, \texttt{deceptive}, \texttt{conceding}, \texttt{unknown}) derived from bid
concentration, trajectory shape, issue-flip rate, and echo detection---not from the opponent's declared agent class.
This avoids brittle, hard-coded routing while still allowing mode-dependent aspiration slopes, closing weights and bait guards.

\section{Agent Architecture}

Agent360 extends the NegMAS \texttt{SAOCallNegotiator} protocol with three cooperating layers.

\begin{enumerate}
  \item \textbf{Bidding persona} --- what we offer: decoy pool construction, phased transition,
        and closing bids.
  \item \textbf{Opponent model} --- what we infer: blended Smith estimators, trajectory features,
        and mode classification.
  \item \textbf{Deal extraction} --- when we accept: time and mode-dependent aspiration, bait
        detection, and deadline safeguards.
\end{enumerate}

On each step with a partner offer, the agent updates all estimators, evaluates acceptance, and otherwise samples the next bid from the current phase.
The first negotiator (seat~0) follows a stricter concealment profile because opening agents
expose more of their bid stream to the opponent's learner.

\section{Concealing Bidding Strategy}

Relative time $t \in [0,1]$ divides the negotiation into three phases.

\paragraph{Decoy phase.}
We bid from a \textbf{maximal-mismatch decoy pool}: rational outcomes that disagree with our
inferred true top-issue values on at least half of the issues, subject to a utility floor.
When we open, transition is blocked until the opponent has made at least four offers
reducing early leakage to curve-fitting learners.
Decoy selection avoids recent repeats, breaks mirror tit-for-tat by perturbing echoed bids, and
prefers offers with large utility jumps relative to our last bid.
The first seat extends the decoy window to $t < 0.40$ (versus $0.35$ for the second seat).

\paragraph{Transition phase ($t < 0.75$).}
We mix decoy outcomes with a widening band of true high-utility outcomes.
The first negotiator keeps decoy mixing for most of the transition window
(\texttt{FIRST\_TRANSITION\_DECOY\_MIX\_UNTIL}${}=0.85$) and reveals true preferences more
slowly (\texttt{FIRST\_TRANSITION\_PROGRESS\_SCALE}${}=0.72$).
If the opponent is classified as conceding and the trajectory shows honest concession, we allow
an \emph{early exit} from decoy even before the usual time gate---trading a small amount of
concealing for advantage when the partner is unlikely to exploit the signal.

\paragraph{Closing phase ($t \ge 0.75$).}
We sample rational outcomes and score each candidate by a blend of normalized self-utility and
estimated opponent utility.
The opponent weight increases through closing but is capped by mode
(e.g.\ higher caps against learners and conceders, lower against deceptive or mirroring
opponents).
Against deceptive late bait-switching, inflated Smith scores are discounted when selecting
closing bids.

Together, these mechanisms implement a \emph{gradual preference revelation} strategy: the
opponent's learner is fed a sustained false early signal, then gradually corrected only when
deal extraction dominates further concealment.

\section{Opponent Modeling}

All opponent inference uses only the observed bid stream, relative time, and negotiation seat.
We never read the opponent's agent class, true utility function, or internal state.
The modeling stack has three roles: estimate what the opponent values, classify how they are
negotiating, and publish a blended utility function that downstream bidding and acceptance
logic can query.

\subsection{Smith-style frequency estimation}

The base estimator is a \textbf{Smith frequency model}: for each issue, we count how often the
opponent chose each value in past offers.
Values repeated often are treated as more valuable to them; per-issue scores are normalized and
averaged into a scalar utility for any outcome.
This matches the learning bias of many league agents (BOA, MAP, MiCRO) and provides a stable
prior over the full negotiation.

We maintain three refinements of this idea:
\begin{itemize}
  \item \textbf{Recency-blended Smith} (second seat): when at least two recent offers exist in
        a window of five bids, we blend full-history Smith with recent-only Smith.
        Weight on the recent estimate grows with sample count (up to ${\approx}0.68$), so we
        track learners that adapt quickly after our persona shifts.
  \item \textbf{Time-weighted Smith}: bids after $t \ge 0.40$ are counted three times when
        building effective frequencies.
        Late offers often reflect genuine priorities after decoy posturing ends; up-weighting
        them reduces contamination from early misleading bids---both theirs and ours.
  \item \textbf{Issue-weighted Smith}: issues where the opponent repeats the same value (low
        spread) receive lower weight; issues with diverse values receive higher weight.
        This filters decoy noise on ``frozen'' issues while emphasizing dimensions where real
        trade-offs appear.
\end{itemize}

\subsection{Trajectory and concession dynamics}

The \textbf{offer trajectory model} records pairs $(t, \hat{u}_{\text{Smith}})$ at each opponent
bid and treats the stream as a time series of revealed preference strength and how much utility the opponent would gain.
From it we compute:
\begin{itemize}
  \item \textbf{Concession slope} --- linear trend of Smith utility over time; slopes below
        ${\approx}-0.04$ indicate honest monotonic concession.
  \item \textbf{Non-monotonicity} --- two or more significant utility \emph{increases} suggest
        strategic feints or bait-and-switch rather than steady concession.
  \item \textbf{Predicted utility at $t$} --- extrapolation used to test whether a late offer is
        inconsistent with the opponent's prior trajectory (a bluff signal).
\end{itemize}
Trajectory features feed both mode classification and bait rejection at acceptance time.

\subsection{Behavioral mode classification}

Rather than routing by opponent \emph{type} (which is unknown in competition), we classify
\emph{behavior} into five modes:
\texttt{mirror}, \texttt{learner}, \texttt{deceptive}, \texttt{conceding}, and \texttt{unknown}.
The decision tree uses mirror echo detection first (at least three matches to our recent bids in
a window of four opponent offers), then trajectory sample count, Smith concentration, concession
slope, and concealment-tactic heuristics.

\paragraph{Deceptive signals.}
We label an opponent deceptive when their bid stream shows concealment tactics, for example:
identical early offers, high early issue-flip rate (${\ge}0.25$), wide early Smith spread
(${\ge}0.12$), non-monotone utility paths, or late bait-switching (preferred values flip on
many issues after $t \ge 0.40$).
Plain BOA/MAP-style learners usually fall into \texttt{learner}, not \texttt{deceptive}, which
limits false bait triggers against honest frequency learners.

\paragraph{Mode-dependent control.}
The closing cap limits how much opponent utility weights closing bids; the aspiration slope
controls how quickly we lower our acceptance threshold over time.
Against deceptive opponents, inflated Smith estimates are discounted before bid scoring and
late offers may be rejected as bait when they diverge from the trajectory prediction.
Conceding opponents may trigger early decoy exit when concession slope is honestly negative;
mirror opponents skip adaptive Smith blends so tit-for-tat echoes do not corrupt the published
model.

\subsection{Published opponent utility blend}

The league measures Concealing from how well the \emph{opponent} models \emph{us}; symmetrically,
we publish \texttt{opponent\_ufun} in private information so the opponent can score Advantage.
The published estimate is built in stages:
\begin{enumerate}
  \item Start from full-history Smith.
  \item If enough late bids exist ($t \ge 0.40$), blend toward a late-phase Smith estimate;
        blend weight increases with late sample count (cap ${\approx}0.55$--$0.68$, higher when
        we open).
  \item If at least three late offers exist, further blend toward issue-weighted late Smith.
  \item On the second seat, blend toward recency Smith when recent data exist; for
        learner/conceding/unknown modes, add a stable-issue match term when recency is strong.
  \item If mode is \texttt{mirror}, skip adaptive blends and return plain Smith.
  \item Clamp all values to $[0,1]$.
\end{enumerate}
Closing bid selection and bait detection query this pipeline; against deceptive opponents,
selected Smith utilities may be discounted before scoring candidate agreements.

\subsection{Anti-mirror response}

When the opponent repeats our last bid, mirror detection fires and we break synchrony by
perturbing the next offer from an anti-mirror pool (small random change on one issue).
Tit-for-tat mirroring freezes preference learning on both sides and depresses
concealing without improving advantage.

\section{Acceptance Strategy}

The agent accepts a partner's offer when any of the following holds:
\begin{itemize}
  \item \textbf{Catastrophe floor:} $t \ge 0.95$ and utility $\ge$ reserved value.
  \item \textbf{Mode aspiration:} utility $\ge u_{\max}(1 - \text{slope} \cdot t)$ and above
        reserved value; the aspiration slope depends on the classified opponent mode.
  \item \textbf{AC-next:} offer utility is at least that of our next concealing bid.
  \item \textbf{Deadline safe:} $t > 0.90$ and utility above reserved value.
\end{itemize}

Before the deadline-safe window ($t \le 0.90$), offers that pass aspiration may still be
rejected if they appear to be \textbf{bait}: the opponent is classified as deceptive, shows
concealment tactics, and the Smith-estimated utility exceeds the trajectory prediction by more
than a threshold (\texttt{ACCEPT\_BAIT\_THRESHOLD}${}=0.14$).

During development, we tested an \emph{escape tier} that accepted above reserved value from
$t \approx 0.88$ onward, intended to reduce timeout penalties.
Although this improved some local stress tests, tournament and proxy runs showed that it
accepted too many mediocre or deceptive agreements against the real opponent mix.
Currently, the agent deliberately disables mid-game escape acceptance and relies on catastrophe accept
only near timeout.

\section{Evaluation and Development Process}

We evaluated each candidate agent on a repeatable local panel before submission:
\begin{itemize}
  \item \textbf{Learners:} BOA, MAP, MiCRO across competition scenarios (both seats).
  \item \textbf{Stress:} time-based NegMAS negotiators (Boulware, Conceder, Tough, etc.).
  \item \textbf{Sparring:} in-house deceptive opponents (mirror, bait-switch, strong learner
        variants).
\end{itemize}
Each matchup cell is repeated with fixed step counts; we report mean Advantage, Concealing, and
their sum.

Table~\ref{tab:local} summarizes a representative development snapshot on the combined panel.
Official tournament runs use a different opponent mix and panel size; we therefore treated
leaderboard feedback as the final gate for acceptance and closing changes, not proxy means
alone.

\begin{table}[ht]
  \centering
  \caption{Local proxy panel (mean Score per matchup cell; development snapshot).}
  \label{tab:local}
  \begin{tabular}{@{}lrrr@{}}
    \toprule
    Opponent group & Advantage & Concealing & Score \\
    \midrule
    Learners (BOA/MAP/MiCRO) & 0.56 & 0.71 & 1.27 \\
    Stress (time-based)      & 0.53 & 0.90 & 1.43 \\
    Deceptive sparring       & 0.61 & 0.58 & 1.19 \\
    \midrule
    Panel mean               & 0.55 & 0.79 & 1.34 \\
    \bottomrule
  \end{tabular}
\end{table}

\paragraph{Iterative refinement.}
The path from an early submission (strong concealing but weak advantage against learners) to
the current Agent360 added: conceding-triggered early decoy exit, relaxed bait-discount thresholds,
mode-dependent closing caps, first-seat decoy rotation, anti-mirror bid pools, and catastrophe
acceptance near timeout.
Changes that improved isolated local cells but hurt the broader distribution---stricter
first-seat gates, shorter second-seat decoy windows, mid-game escape acceptance, and
stall-based accept rules---were removed after broader evaluation showed regressions.

\section{Discussion and Conclusions}

Agent360 is a concealment-first bilateral negotiator grounded in signaling theory
and empirical tests of alternative personas.
Its core idea is psychologically motivated but implemented entirely through bid streams:
early maximal-mismatch decoys mislead frequency learners, gradual transition limits preference
leakage, and mode-aware closing and acceptance extract Advantage without surrendering
Concealing to deceptive or mirroring opponents.

Three lessons emerged from development.
First, \textbf{concealing and advantage require separate controls}; strengthening acceptance
without preserving the decoy persona hurts the combined score against learner-heavy fields.
Second, \textbf{evaluation must span multiple opponent families}; learner, stress, and deceptive
sparring panels together catch failures that no single benchmark family reveals alone.
Third, \textbf{seat asymmetry matters}; first and second negotiators see different information,
and asymmetric decoy length and min-offer gates improved learner matchups without oracle
routing.

The submitted agent operates without privileged information, adapts to inferred opponent behavior,
and balances strategic misrepresentation of preferences with profitable agreement formation.

\end{document}
